\documentclass{ctexbook}
\usepackage{amsmath,amssymb,amsthm}
\usepackage[hidelinks]{hyperref}
\newcommand{\corr}[1]{\footnote{本批译校：#1}}
\newcommand{\ednote}[1]{\footnote{英译编者注：#1}}
\newcommand{\Ys}{Y^{\mathrm{sing}}}
\newcommand{\Yss}{Y^{\mathrm{sup\,sing}}}
\title{代数几何基础第五卷预备笔记\\8.9至8.11：校订续译}
\author{据用户上传英译编辑稿译出}
\date{2026年11月14日}
\begin{document}
\maketitle
\frontmatter
\chapter*{范围与未解决的前提}
本批从上传底本 \texttt{EGAVaiello.pdf} 第29页“Corollary 8.9.”开始，连续译至第30页（8.11）证明末，止于“Remark 8.12.”之前。包括（8.9）、其证明及关于（8.8）的收尾说明、（8.10）、（8.11）和原脚注41--43。

底本为 Piotr Blass 与 Joseph Blass 翻译编辑的 Grothendieck 第五卷预备笔记，未与法文手稿核对。（8.7）(c)的切空间结论及（8.8）的标题、完整假设仍缺失。上一批已证实：在现存“$f$非分歧”假设下，泛平展不能推出双有理，因而（8.9）的开浸入断言也不能直接使用。本批\textbf{不猜补缺失条件，不把原稿未成立的证明改称已完成证明}。

按用户要求，已核实的正确结论写入正文；已证实的错误与原推演在脚注中完整译录，并说明理由。尚不能确定的片段标为待核。局部源靶和符号的确定修正也在脚注留痕，独立补证与原著文字分开。

统一记号：$Q=\Ys$，$R=\Yss$，$q:Q\to P^\vee$；$T$ 为像闭包的约化结构，$g:Q\to T$，$U=T_{\mathrm{sm}/k}$。在使用函数域时假设 $Q,T$ 非空且整，并沿用有限型背景；$K=k(T)$、$L=k(Q)$，$r=\dim P^\vee$，$\dim Q=r-1$。$R$ 按（8.7）(b)为 $q$ 的分歧闭集。$Q_U=g^{-1}(U)$。这些条件在下文使用时明示，不声称已恢复原稿缺失的整套假设。

查阅原对话及Library最新独立译稿，未发现本范围已有交付；原会话仍为第七版918页、第0章至（19.6.7），独立第五卷恰止于（8.9）标题前。本批不覆盖旧工程。
\mainmatter
\chapter*{8.9\quad 推论：现存条件下的可核部分}
假设（8.8）的条件(i)成立，即 $g$ 泛平展，令 $U$ 为 $T$ 中相对于 $k$ 光滑的点所成的开集。可以确认 $U\ne\varnothing$。至于 $Q_U\to U$ 是开浸入的结论，在现存假设下不成立；原稿完整断言及固有情形的后续结论见脚注，均不能仅据泛平展性采纳。\corr{原推论译文：“假设（8.8）的条件(i)成立，令 $U$ 为 $T$ 中相对于 $k$ 光滑的点所成的开集。则 $U\ne0$，$\Ys/U\to U$ 为开浸入，因而 $\Ys/U$ 不含 $\Yss$ 的点。若 $X$ 在 $K$ 上固有，则 $g:\Ys\to T$ 满射，从而 $\Ys/U\to U$ 在 $K$ 上固有，所以 $g:\Ys\to T$ 满射，因此 $\Ys/U\to U$ 是同构，且 $U$ 是具有后一性质的最大开集。”这里两个 $K$ 与前文函数域记号冲突；若讨论 $X$ 的基域，应为 $k$。此外，$U$ 上的态射固有应指相对于 $U$，不能推出其源在 $k$ 上固有。更重要的是，取 $k=\mathbb C$、$X=\mathbb G_m$、$f(t)=[1:t^2:t^4]$，则 $f$ 非分歧，$Q\simeq X$，$g(t)=[t^4:-2t^2:1]$ 到光滑对偶圆锥曲线泛平展、次数二；此时 $U=T$，$Q_U\to U$ 不是开浸入。故本批不把这条主结论作为正确正文。反例只针对现存假设，不能否定未知完整假设下的原定理。}

如果另已证明 $Q_U\to U$ 为开浸入，并且该态射固有且满射，那么它必为同构。这是原稿后半段可以单独核实的条件式，并不填补其前半段的缺口。\corr{固有的开浸入，其像同时开且闭；再加满射便是同构。此论证没有证明当前 $g$ 满足开浸入假设。若 $Q$ 在 $k$ 上固有而 $T$ 在 $k$ 上分离，则 $g$ 固有；其支配性与固有性使像为全 $T$。这是正确的固有性传递方向。}

\section*{证明与缺口记录}
首先，由 $g$ 支配且泛平展，可以找到 $T$ 的一个非空开子集 $V$，使 $g^{-1}(V)\to V$ 平展且满射；于是 $V$ 在 $k$ 上光滑，从而 $U$ 非空。\corr{原文把“泛平展”重复写了两次，并混用 $k$ 与 $K$；此处只去掉重复，统一基域为 $k$。这里的 $Q$ 在 $k$ 上光滑，光滑性沿平展满覆盖下降；泛平展且泛有限使适当缩小 $T$ 后可除去非平展部分的像，再利用平展像开性得到满射。}

对于 $\xi\in V(k)$ 及其上方的 $y\in Q(k)$，原稿试图利用（8.7）(c)从 $T$ 在 $\xi$ 处的切空间唯一恢复 $x$，进而恢复 $y$。这一唯一性推断没有得到现存假设的支持。\corr{原证明对应段落完整译为：“若 $\xi\in V(K)$，且 $y$ 是其上方 $\Ys(k)$ 的一点，则沿用（8.7）c)的记号，空间 $H^0$ 正是 $T$ 在 $\xi$ 处的切空间；如我们在那里所指出的，$y$ 投影到 $X(k)$ 的点 $x$ 由与 $H^0$ 正交的点确定，因而唯一确定；又因为 $\Ys\subset X\times P^\vee$，所以 $y$ 也唯一确定。这已证明 $g$ 是双有理的，因为它泛平展且泛根性。”但(c)在底本中残缺，$H^0$ 也未完整定义。即使切空间确定了环境射影空间中的 $f(x)$，一般非分歧态射仍可能在此点上方有多个 $x$；二次覆盖反例正是如此。故不能从这段推断双有理性。$K$ 的大小写和 $H^0/H^1$ 差异均是底本原状。}

接下来构造逆映射的段落同样依赖该唯一性；因此本批保存其完整可辨文字作为待核原证明，不把它改造为已经建立的逆态射。\corr{原文续译：“另一方面，考虑态射 $\psi$（其定义形式显然）：它给每个 $\xi\in U(K)^0$ 对应 $P$ 中与 $U$ 在 $\xi$ 处切空间正交的唯一点 $x=\psi(\xi)$；在 $V$ 上，它与复合 $V\to\Ys|_V\to X$ 一致，第二个箭头为投影。因此，令 $h=(\psi,\mathrm{id}):U\to P\times T$〔字迹不清〕，$g_1=g|_{g^{-1}(U)}:g^{-1}(U)\to U$，则复合 $hg_1:g^{-1}(U)\to P\times T$ 恰为典范包含，因为在 $g^{-1}(V)\xrightarrow{\sim}V$ 上它是如此。于是 $h$ 经由 $P\times T$ 中 $Y^1$ 的概形论闭包 $Y^1$ 分解；所以 $Y^1$（它在上述闭包中开）在 $h$ 下的逆像是 $U$ 的开子集，记为 $U^1$。由于 $hg_1=\text{包含}$，立即可见 $g^1$ 诱导同构 $g^{-1}(u)\xrightarrow{\sim}W$ 是同构；由 $Y^1$ 光滑可知 $W$ 光滑，所以 $W\subset U$。这就证明了（8.9）。”底本在闭包与原子概形处均印作 $Y^1$，并有 $U^1,W,u,g^1$ 混用；这些对象尚不能唯一辨认，未据猜测修复。原脚注41为“请询问 A.G.”，已在此保留。即便处理记号混乱，此段仍需要前面未证明的唯一性与浸入类条件。}

\section*{关于（8.8）末尾维数断言的补充}
若 $Q$ 非空、整且维数为 $r-1$，$R$ 为真闭子集，则 $R$ 为空，或者 $\dim R\leq r-2$。要把不等式加强为非空时恰等于 $r-2$，还须证明这里的 $R$ 确由一个非零可逆模截面定义；本批不据缺字补出这一识别。\corr{原文：“（8.8）最后的断言 $\Yss=\varnothing$ 或 $\dim\Yss=r-2$，以及 $\dim T=r-1$，都是平凡的：第一个来自 $\Ys$ 不可约且维数为 $r$，以及 $\Ys$ 或 $\Yss$〔字迹不清〕是可逆模某截面 $D$ 的消失所定义的这一事实。”其中 $\dim\Ys=r$ 与前面射影法丛的计算及本句结尾均冲突，应为 $r-1$。底本原脚注42为“请询问 A.G.”。本批不将有待辨认的零点对象和截面认定为已经恢复。}
第二个维数等式来自 $L/K$ 有限，因而
\[
\operatorname{trdeg}_k L=\operatorname{trdeg}_k K,
\qquad\dim T=\dim Q=r-1.
\]

\chapter*{8.10\quad 备注与集合差的订正}
沿用（8.9）的记号，在已另行获得该推论所需前提和开浸入结论时，有
\[
g^{-1}(U)\subset Q-R.
\]
原稿指出，即使 $X$ 闭于 $P$，这一包含也未必是等号。令 $E_{\mathrm{et}}$ 为 $g$ 的平展点集，则可独立确定：“非分歧而不平展”的点集是
\[
(Q-R)\setminus E_{\mathrm{et}}.
\]
只有另已证明 $E_{\mathrm{et}}=g^{-1}(U)$ 时，才能把上式写成 $(Q-R)\setminus g^{-1}(U)$。\corr{底本说“$\Ys-\Yss$ 是 $g$ 非分歧但不平展的点集”，漏掉了集合差。按（8.7）(b)及 $T\hookrightarrow P^\vee$ 为闭浸入，$Q-R$ 是 $g$ 的非分歧点集，其中还包含平展点，必须去掉平展点集。原稿把 $g^{-1}(U)$ 认作平展域，仍依赖待核的（8.9）背景，本批不将两者直接等同。}

在本段有限型、整、支配、同维的背景下，非分歧而非平展点的像不是几何单支点，因而也不是几何正规点。\corr{原文为“这还蕴含 $g(y)^0$ 是 $T$ 中一个不几何正规、甚至不几何单支的点”；上标 $0$ 含义不清，正文只写像点 $g(y)$。独立核验的理由是：非分歧态射在靶的平展局部可写成闭浸入到一个平展邻域；若靶在该点几何单支，则相应平展局部的约化结构只有一个局部分支。在这里靶整、局部平展代数约化，因而相应局部环整。缩小到该点的整平展邻域后，其维数与源均等于 $\dim T$；闭浸入的像因同维而占满这个整邻域，故其定义理想为零，故此局部闭浸入为同构，原态射在该点平展，矛盾。这个结论使用了本段的支配、同维条件，不能仅凭“非分歧但不平展”推出，也不证明（8.9）的开浸入断言。}

从几何上说，这对应于如下现象：一个超平面在 $x\in X(k)$ 处与 $X$ 相切但不密切，同时还可能存在另一点 $x^1\in X(k)$，使同一超平面在 $x^1$ 处相切。\corr{底本句末写“at $x$ [or $x^1$ illegible]”，译为“在 $x$ 处〔或者是 $x^1$，字迹不清〕”；由该句明确讨论另一接触点的语义，正文用 $x^1$。原脚注43为“字迹不清”，保留于此。不密切仍按前文解释为相应二次型非退化，而非随意替换为其他相切阶数定义。}
原稿举例说，可以令 $P$ 维数为二、$X$ 为次数至少四的非奇异曲线，并称在任意特征下都有显然的例子。〔这里是在转述原稿的例子预告，本批未把它加强为“所有这样的曲线都有双切线”。〕〔附注：$X$ 的“双切线”对应于“对偶曲线”的双点。〕

\chapter*{8.11\quad 特征零的推论}
假设 $k$ 的特征为零，保持上文 $Q$ 非空且整、光滑、维数 $r-1$ 等背景。

\noindent\textbf{(a)} $R=\Yss$ 在 $P^\vee$ 中的像的闭包余维至少为二。

\noindent\textbf{(b)} 条件(iii)，即 $\dim T=r-1$，等价于 $L/K$ 有限，亦等价于(i bis)、(i)及(ii)。这里(ii)为 $Q\ne R$。因而其否定 $Q=R$ 等价于 $Q$ 在 $P^\vee$ 中的像的闭包余维至少为二。\corr{原文说“（8.8）的条件(iii)与其他条件等价”，继而说明其他条件的否定，即 $\Ys=\Yss$，意味着 $\Ys$〔或 $\Yss$，字迹不清〕在 $P^\vee$ 中的像余维至少二。若“其他条件”包含(i ter)的双有理性，则这一说法即便在特征零也被二次平展覆盖反例否定；若包含(ii bis)的 $k$-有理点存在，还缺少代数闭等条件。本批只保留可由现存整性、维数、有限型和特征零条件证明的等价关系。最后的难辨对象在 $Q=R$ 时无底集差别，故正文写 $Q$，并在此留存原文两种读法。}

\begin{proof}
考虑到前述维数结论，(b)蕴含(a)。由维数理论，(iii)意味着 $L/K$ 为有限扩张；也可以把它写成（8.8）中的条件“(iii bis)”。而特征零时 $L/K$ 总是可分的，所以得到（8.8）的(i bis)。这就是原稿给出的论证。

为避免依赖缺失的全组等价断言，补充说明如下〔本批补证〕：$T$ 是 $q$ 的像闭包，故 $\dim T\leq r-1$。若 $\dim T\leq r-2$，则 $R$ 的像闭包包含于 $T$，(a)立即成立；且 $q$ 不可能在任何点非分歧，否则整性使其在泛点非分歧，从而 $L/K$ 有限，迫使两者维数相等。故此时 $R=Q$。

若 $\dim T=r-1$，则 $L/K$ 有限，特征零保证可分，因此 $g$ 泛平展，$q$ 泛非分歧。由（8.7）(b)，$R$ 为 $Q$ 的真闭子集，于是 $R$ 为空或维数至多为 $r-2$，其像的闭包余维至少为二。反之，$Q\ne R$ 使非分歧开集非空并含泛点，故 $L/K$ 有限可分，得到(iii)。这证明了正文中明确列出的等价关系，也直接证明(a)，没有使用未经证实的双有理性。
\end{proof}
\backmatter
\chapter*{核查与接续}
本批完整覆盖（8.9）至（8.11）在所上传英译编辑稿中的连续文字。不能成立或不能辨认的证明步骤已在脚注中完整译录，不假造逆映射、不恢复缺失的（8.8）。正文独立保留可确认的非空光滑开集、正确的集合差和特征零余维结论；附加证明标明为本批补证。

二次覆盖反例仅否定现存“非分歧”假设下的开浸入、双有理推断，不否定缺失的完整假设下可能成立的原结论。（8.10）的几何单支结论按显写的整性、支配及同维条件另作核验，原稿任意特征例子的细节仍待核。旧批次未在本轮整体回溯修订，亦未合并原合订本。

\textbf{下一断点：}底本第30页“Remark 8.12.”标题。下次应从该备注开始，继续处理其跨至第31页的段落，不要跳到（8.13）。
\end{document}
